\chapter[ECP5 implementation \& benchmarks]{ECP5 implementation, real benchmark campaign \& measured results}
\label{ch:impl2}

\section{Flow and verification discipline}
Every number in this chapter is labelled \textsc{Theoretical},
\textsc{Simulated}, \textsc{Post-P\&R measured}, or \textsc{Derived}
(a combination of two real measurements, e.g.\ cycles $\div$ real Fmax).
No result is invented, approximated to look better, or reported without
a matching real measurement.

\begin{tabularx}{\textwidth}{L{5.0cm} C{3.0cm} Y}
\toprule
\rowh \thd{Verification stage} & \thd{Outcome} & \thd{Covers} \\
\midrule
RTL simulation (Verilator 5.050) & \PASS & bit-exact correctness vs.\ a software golden model \\
\rowa ECP5 synthesis (Yosys \code{synth\_ecp5}) & \PASS, 0 problems & synthesizability, resource mapping \\
Place\&route (real \code{nextpnr-ecp5}) & \PASS at \code{N\_SLOTS}$\le$2 & LUT/FF/DSP, real timing \\
\rowa Full benchmark campaign & 24/24 bit-exact & 6 workloads $\times$ 4 \code{N\_SLOTS} configurations \\
\bottomrule
\end{tabularx}

\begin{fnnote}[Verilator, not Icarus, for V2]
Two independent Icarus Verilog v13.0 scheduling defects were found and
reproduced on minimal repros during V2's own M1 milestone (a
task/scope-entry desync and a spurious condition evaluation, both
edge-parity dependent) --- Verilator gives correct results on the same
repros. V1's own certification (performed separately, with Icarus) was
unaffected, since its own testbenches already avoided the trigger
pattern by convention; this is flagged honestly, not glossed over.
\end{fnnote}

\section{V1 vs.\ V2 --- final comparison}
Both systems full-system (not isolated modules), same
PARALLEL/P\_IN=8, same real, unmodified V1 PSRAM chain.

\begin{tabularx}{\textwidth}{L{3.6cm} C{2.6cm} C{2.6cm} Y}
\toprule
\rowh \thd{Metric} & \thd{V1} & \thd{V2 (N\_SLOTS=2)} & \thd{Class} \\
\midrule
Fmax & 68.65~MHz (\FAIL) & \textbf{87.72~MHz} (\PASS) & Post-P\&R \\
\rowa LUT (Total LUT4s) & 8907 & 4359 & Post-P\&R \\
FF (Total DFFs) & 4900 & 3924 & Post-P\&R \\
\rowa DSP (MULT18X18D) & 16 & 16 & Post-P\&R \\
BRAM (DP16KD) & 2 & 0 & Post-P\&R \\
\rowa cycles/neuron (1 neuron, 8 inputs, real PSRAM) & 209 & 166 & Simulated \\
Real wall-clock speedup vs.\ V1 & 1.00$\times$ & \textbf{2.6$\times$} & Derived \\
\bottomrule
\end{tabularx}
\begin{center}\footnotesize\itshape\color{fnGrey}
V1's own figures are its already-certified, frozen baseline (not
re-measured this session); V2's figures are real, current measurements
including both post-campaign optimizations.\end{center}

\begin{fnnote}[Where the win comes from --- and where it does not]
V2's advantage comes from a faster pipeline and a higher achievable
clock, \textbf{not} primarily from the multi-processor concurrency the
architecture was built to add. That concurrency's own real payoff, given
the single-PSRAM-port memory subsystem, is much smaller than a naive
\code{N\_SLOTS}$\times$\code{P\_IN} calculation would suggest ---
\S\ref{sec:scaling}.
\end{fnnote}

\section{Real \texttt{N\_SLOTS} sweep --- Fmax and resources}
Full system, real place\&route, both memory optimizations active
(word-burst \S\ref{sec:burstimpl} $+$ activation cache
\S\ref{sec:cacheimpl}).

\begin{tabularx}{\textwidth}{C{1.6cm} C{2.2cm} C{1.6cm} C{1.6cm} C{1.6cm} C{1.6cm} Y}
\toprule
\rowh \thd{N\_SLOTS} & \thd{Fmax} & \thd{LUT4} & \thd{FF} & \thd{DSP} & \thd{BRAM} & \thd{80\,MHz} \\
\midrule
1 & 131.79~MHz & 2760 & 2405 & 8/72 & 0 & \PASS \\
\rowa 2 & \textbf{87.72~MHz} & 4359 & 3924 & 16/72 & 0 & \PASS (\textbf{recommended}) \\
4 & 65.01~MHz & 9158 & 7986 & 32/72 & 0 & \FAIL \\
\bottomrule
\end{tabularx}

\subsection{Fmax versus N\_SLOTS}
\begin{center}
\begin{tikzpicture}
\begin{axis}[
  width=0.68\textwidth,height=6.0cm,
  xlabel={\footnotesize N\_SLOTS}, ylabel={\footnotesize Fmax [MHz]},
  xtick={1,2,4}, ymin=50,ymax=145, ytick={50,65,80,95,110,125,140},
  tick label style={font=\scriptsize}, label style={font=\footnotesize},
  grid=major, grid style={fnRule!40},
  legend style={font=\scriptsize,at={(0.5,-0.28)},anchor=north,legend columns=2}]
 \addplot[fnBlue,mark=square*,thick,mark options={fill=fnBlue}]
   coordinates {(1,152.44)(2,133.58)(4,112.07)};
 \addlegendentry{word-burst only}
 \addplot[fnRed,mark=*,thick,mark options={fill=fnRed}]
   coordinates {(1,131.79)(2,87.72)(4,65.01)};
 \addlegendentry{$+$ activation cache}
 \draw[fnAmber,dashed,thick] (axis cs:1,80)--(axis cs:4,80);
 \node[font=\scriptsize,text=fnAmber] at (axis cs:3.3,74){80 MHz target};
\end{axis}
\end{tikzpicture}
\end{center}
\begin{center}\footnotesize\itshape\color{fnGrey}
The activation cache's own real Fmax cost grows much faster with
N\_SLOTS than the arbiter-widening cost alone --- a single shared
resource with N\_SLOTS request ports and an unpipelined,
broadcast-capable hit-check.\end{center}

\section{Real parallel scaling}
\label{sec:scaling}
Not assumed --- computed from real cycle counts, largest workload
(256 independent neurons sharing one input vector).

\begin{tabularx}{\textwidth}{C{1.6cm} C{2.6cm} C{2.0cm} C{2.6cm} C{2.0cm}}
\toprule
\rowh \thd{N\_SLOTS} & \thd{Speedup(N)} & \thd{Efficiency} & \thd{PSRAM utilization} & \thd{Real wall-clock speedup vs.\ N=1} \\
\midrule
1 & 1.00$\times$ & 100\% & 55.5--71.8\% & 1.00$\times$ \\
\rowa 2 & 1.06$\times$ & 53\% & $\approx$90\% & 0.99$\times$ (a wash) \\
4 & 1.06$\times$ & 27\% & $\approx$90\% & 0.79$\times$ (\emph{slower}) \\
8 & 1.06$\times$ & 13\% & $\approx$90\% & --- \\
\bottomrule
\end{tabularx}
\begin{center}\footnotesize\itshape\color{fnGrey}
Pre-optimization figures, isolating the real scaling behavior from the
two memory optimizations' own effect (\S\ref{sec:burstimpl}--\ref{sec:cacheimpl}).\end{center}

\begin{fnwarn}[The central, measured finding]
Real cycle-count speedup from \code{N\_SLOTS}=1 to \code{N\_SLOTS}=8 is
essentially flat (1.05--1.06$\times$) for sustained, memory-bound
workloads --- the single shared PSRAM port saturates at
$\approx$90\% utilization regardless of \code{N\_SLOTS}$\ge$2. Once real
Fmax degradation is also folded in, \code{N\_SLOTS}=4 measures
\emph{slower} in real wall-clock time than \code{N\_SLOTS}=1. More
hardware parallelism made this workload class worse, not better,
because the bottleneck was never compute.
\end{fnwarn}

\section{Memory optimization \#1 --- word-level burst reads}
\label{sec:burstimpl}
See ch.~\ref{ch:mem}, \S\ref{sec:burst}, for the full rationale. Real,
measured single-job cycle reduction: $-$49\% (1 tile), $-$54\% (3
tiles), $-$56\% (5 tiles). Combined real wall-clock speedup on the
256-neuron sustained workload: 2.24--2.37$\times$ across every
\code{N\_SLOTS} tested, at negligible real Fmax cost.

\section{Memory optimization \#2 --- shared activation cache}
\label{sec:cacheimpl}
See ch.~\ref{ch:mem}, \S\ref{sec:cache}. A further real 1.66--2.00$\times$
cycle reduction on top of optimization~\#1, at a real, steep Fmax cost
that makes \code{N\_SLOTS}=4 fail 80\,MHz outright.

\begin{tabularx}{\textwidth}{C{1.6cm} C{2.4cm} C{2.4cm} C{2.4cm}}
\toprule
\rowh \thd{N\_SLOTS} & \thd{Wall-clock, baseline} & \thd{Wall-clock, final} & \thd{Total real speedup} \\
\midrule
1 & 5118.1~$\mu$s & 1324.9~$\mu$s & \textbf{3.86$\times$} \\
\rowa 2 & 5169.5~$\mu$s & 2113.9~$\mu$s & \textbf{2.45$\times$} (recommended) \\
4 & 6498.7~$\mu$s & 2842.6~$\mu$s (Fmax fails) & 2.29$\times$ but a real regression vs.\ \#1 alone \\
\bottomrule
\end{tabularx}

\section{Bottleneck analysis}
\begin{tabularx}{\textwidth}{L{2.8cm} Y}
\toprule
\rowh \thd{Candidate} & \thd{Verdict, with real evidence} \\
\midrule
\textbf{Memory (PSRAM port)} & \textbf{The real bottleneck.} $\approx$90\% utilization at N\_SLOTS$\ge$2; real compute-to-memory-wait ratio on the order of 1:170--1:220. \\
\rowa Compute (Neural Processor) & Not the bottleneck --- the pipeline is idle most of the time waiting for data. \\
Arbiter overhead & Real but small: a mandatory 1-cycle pending-latch (correctness, not choice) plus modest Fmax cost ($-$6\% at N\_SLOTS=2, word-burst alone). \\
\rowa Director/dependency logic & Not the bottleneck --- zero lost/duplicated jobs, no queueing backlog observed; a real \emph{fairness} issue exists (\S\ref{ch:sched}) but does not limit throughput. \\
DSP/LUT/FF availability & Not the bottleneck at N\_SLOTS$\le$4 --- all well under budget; DSP would eventually bind at N\_SLOTS=9 (64/72), never reached in practice since the memory bottleneck dominates first. \\
\bottomrule
\end{tabularx}

\section{Limitations, honestly stated (PSRAM-era campaign above)}
\begin{itemize}
\item V1's own memory-utilization/stall figures were not re-measured
      this session (V1 is frozen); only its already-certified numbers
      are used for comparison.
\item No clean per-cycle split between ``processor computing'' and
      ``processor waiting for memory'' exists in the current
      instrumentation --- reported figures use tile-delivery-rate
      proxies, not an exact split.
\item Power/energy: \textbf{NOT MEASURED} --- no ECP5 power estimator
      (\code{ecppower}, \code{icepower}, or equivalent) is available in
      this project's toolchain; no value is invented in its place.
\item \code{N\_SLOTS}=8 was not re-measured full-system (real PSRAM
      chain) after either memory optimization --- only
      \code{dataflow\_core.v} alone, pre-optimization (92.63~MHz).
\end{itemize}

\section{SDRAM-era benchmark addendum (2026-09-07) --- current,
authoritative results}
\label{sec:impl-sdram-addendum}
\begin{fnwarn}[Supersedes the PSRAM/\code{N\_SLOTS}$\le$2-era campaign
above for the current hardware baseline]
Every section above (V1 vs.\ V2 comparison, \code{N\_SLOTS} sweep,
parallel scaling, memory optimizations \#1/\#2, bottleneck analysis)
describes an earlier V2 milestone built on V1's own PSRAM chain,
recommending \code{N\_SLOTS}=2. The project has since replaced external
memory with a single SDR SDRAM device (ch.~\ref{ch:mem}
\S\ref{sec:sdram-mem-addendum}) and closed on
\textbf{\code{N\_SLOTS}=4 as the production configuration}. This
section is the current, real, measured state; the PSRAM-era numbers
above remain real and correctly measured for the architecture they
describe, but do not apply to the current board.
\end{fnwarn}

\subsection{Real resource utilization (\code{N\_SLOTS}=4, SDRAM
architecture, post real critical-path fixes)}
\begin{tabularx}{\textwidth}{L{4.2cm} C{2.4cm} Y}
\toprule
\rowh \thd{Resource} & \thd{Count} & \thd{Notes} \\
\midrule
TRELLIS\_COMB (LUT4-equiv) & 7,175 / 43,848 (16.4\%) & Real Yosys synthesis, most recent measurement (post-ERR-0029) \\
\rowa MULT18X18D & 32 / 72 (44.4\%) & Exactly $4\times8$ (\code{N\_SLOTS}$\times$\code{P\_IN}), confirmed --- the ERR-0027 fix removed a spurious 33rd multiplier \\
DP16KD (block RAM) & 0 / 108 & All small SRAMs synthesize to distributed RAM \\
\rowa EHXPLLL & 1 & Real \code{EHXPLLL} primitive, \code{ecppll}-derived parameters \\
TRELLIS\_FF & $\ge$6,322 (last individually re-quoted figure) & Real, same SDRAM architecture, pre-dates the ERR-0027/0028/0029 restructuring; not independently re-synthesized standalone since --- disclosed as a lower-bound reference, not re-invented as exact \\
\bottomrule
\end{tabularx}

\subsection{Real, current clock closure and functional regression}
See ch.~\ref{ch:hw} \S\ref{sec:clock-closure-current} for the complete
per-seed Fmax/WNS table (single source of truth, not duplicated here):
\textbf{\code{N\_SLOTS}=4 @ 64\,MHz, 8/8 seeds PASS} (worst 64.55\,MHz,
best 72.37\,MHz); \code{N\_SLOTS}=8 deferred (3/8); 80\,MHz confirmed
NO-GO at either processor count with a genuinely regenerated PLL.

D-Stress functional regression (256 neurons, 256/256 bit-exact vs.\
golden model): \textbf{49,927 cycles} at \code{N\_SLOTS}=4 ---
\textbf{780\,\textmu s} real wall-clock at the P\&R-verified 64\,MHz
system clock ($49{,}927 / 64{,}000{,}000$, \textsc{Derived}). SDRAM
directed boundary verification (ch.~\ref{ch:hw}
\S\ref{sec:sdram-addendum}): 21/21 PASS, zero bugs found, both 64\,MHz
and 166\,MHz.

\subsection{Real SPI host protocol throughput}
Board-level smoke test (\code{tb\_fpga\_neural\_v2\_top\_smoke.v}, 11/11
PASS): single job 99--100 cycles/job; back-to-back 88--100 cycles/job;
steady-state throughput unaffected by inter-job gap (100\,ns/5\,\textmu
s/50\,\textmu s tested). Maximum verified SPI host clock: \textbf{12\,MHz
recommended} (exact deterministic CDC edge at 12.8\,MHz $=$ 64\,MHz/5)
--- see ch.~\ref{ch:hw} \S\ref{sec:spi-max-verified} for the full sweep.

\subsection{Limitations, honestly stated (current SDRAM architecture)}
\begin{itemize}
\item Power/energy: \textbf{NOT MEASURED} --- no ECP5 power estimator
      is available in this toolchain (unchanged from the PSRAM-era
      disclosure above).
\item Hold-time closure: \textbf{OPEN --- tool-chain limitation}, not a
      real defect; see ch.~\ref{ch:hw} \S\ref{sec:hw-open-items} for
      the complete, consolidated open-items list.
\item \code{N\_SLOTS}=8 is functionally correct but not
      timing-closed on every tested seed --- deferred by explicit
      project direction, not attempted further this pass.
\item No embedded-host (ESP32-class) physical baseline exists; all
      host-side numbers above are protocol-level simulation, not
      measured on real silicon.
\end{itemize}
